\chapter{Methodology}

\section{Overview of Technical Approach}
\label{sec:method_overview}

TechPulse has been implemented as a cloud-native Hybrid Retrieval-Augmented Generation system designed to synthesize emerging technology intelligence from heterogeneous sources. The technical approach is organized around five interconnected components: (1) continuous heterogeneous data ingestion from five sources; (2) deterministic preprocessing and semantic embedding; (3) hybrid multi-stage retrieval; (4) context-bound grounded generation with automatic LLM fallback; and (5) production-grade MLOps deployment and monitoring.

The core hypotheses under evaluation are stated formally as:

\begin{itemize}[leftmargin=0.5in]
  \item \textbf{H\textsubscript{0}}: There is no statistically significant difference in mean context precision between hybrid and baseline (vector-only) retrieval, evaluated over 50 matched queries under identical corpus and embedding conditions.
  \item \textbf{H\textsubscript{1}}: Hybrid retrieval achieves statistically higher mean context precision than baseline retrieval (one-sided Wilcoxon signed-rank test, $\alpha = 0.05$).
\end{itemize}

\noindent Secondary evaluation metrics (faithfulness, answer relevancy, and composite score) are analyzed descriptively and are not subject to formal hypothesis testing, serving to characterize the broader impact of hybrid retrieval beyond context precision alone.

A secondary performance hypothesis requires that the hybrid retrieval pipeline does not exceed the p95 latency bound ($\tau = 2.0$~seconds for the retrieval stage; end-to-end latency including LLM generation may exceed this bound and is reported separately as an informational metric) or the monthly cost ceiling (USD~15 at an assumed workload of 1,500 queries/month, i.e., USD~0.01 per query). The evaluation is designed to be informative regardless of outcome: if H\textsubscript{0} cannot be rejected, the study characterises the conditions under which hybrid and baseline converge---an empirically valuable finding given the widespread adoption of hybrid retrieval in production RAG systems without controlled comparative evidence. These hypotheses are tested through a controlled experiment comparing both configurations under identical corpus and embedding conditions.

\section{Why This Approach is Suitable}
\label{sec:method_rationale}

Three considerations motivate the hybrid RAG approach. First, emerging technology knowledge is temporally sensitive; vector-only search lacks recency signals, whereas BM25+RRF fusion explicitly addresses temporal drift. Second, the system must operate within AWS Free Tier cost and latency bounds: full-corpus BM25 adds $O(N)$ computation against a cached index, RRF merges ranked lists in $O(k)$, and a lightweight cross-encoder (MiniLM-L-6-v2, 22M parameters) reranks the small top-$k$ candidate set with minimal added latency. Third, AWS-managed services (Lambda, RDS, API Gateway, EventBridge) enable a student team to deploy a production-grade system within the course timeline and budget.

\section{Technical Contributions and Design Decisions}
\label{sec:innovation}

The engineering contribution of TechPulse lies not in any single component but in their integration into a coherent, cost-governed MLOps system. Specifically, four design decisions represent well-motivated engineering choices that distinguish this system from typical RAG prototypes.

\textbf{True parallel hybrid retrieval with weighted fusion and cross-encoder reranking:} The system performs two independent retrieval passes---cosine vector search (pgvector) and BM25 keyword search (BM25Okapi over the full indexed corpus)---and fuses their rankings with weighted Reciprocal Rank Fusion (RRF; Cormack et al., 2009), where each signal receives a data-informed weight (vector~0.50, BM25~0.35, recency~0.15). A third recency signal is integrated as an additional RRF ranking. Following fusion, a lightweight cross-encoder (ms-marco-MiniLM-L-6-v2) reranks the top candidates to capture fine-grained query--passage interactions that embedding similarity alone cannot model (Nogueira \& Cho, 2019). A design-phase comparison against weighted $\alpha$/$\beta$/$\gamma$ linear reranking showed that linear scoring is sensitive to corpus composition and incurs prohibitive tail latency. In full evaluation, the linear configuration reached p95\,=\,25\,s, violating the retrieval-stage latency constraint ($\tau = 2.0$\,s), and was therefore not suitable under the latency constraints of this system; the empirical evidence is documented in Section~\ref{sec:sensitivity}.

\textbf{Unified vector--relational retrieval without dedicated infrastructure:} PostgreSQL with pgvector combines structured SQL metadata filtering and cosine vector search in a single managed database, eliminating the need for separate vector store infrastructure (Section~\ref{sec:storage}).

\textbf{Idempotent ingestion pipeline with SHA-256 content-addressed deduplication:} Canonicalized document fields are hashed using SHA-256, enabling safe re-runs without duplicate records---supporting experiment replay, crash recovery, and corpus versioning without a separate deduplication service.

\textbf{Budget-aware inference with programmatic halt:} AWS Cost Explorer queries are integrated into the inference path, halting invocations programmatically when monthly cost thresholds are approached. This transforms cost from a passive monitoring concern into an active runtime constraint.

\subsection{Scope of Machine Learning Development}
\label{sec:ml_scope}

TechPulse does not train a foundation model from scratch. Instead, the machine learning contribution lies in the selection, integration, configuration, and controlled evaluation of pretrained components within a production-grade RAG pipeline. The system uses fastembed (ONNX Runtime) with the pretrained embedding model \texttt{sentence-transformers/all-MiniLM-L6-v2}, a pretrained cross-encoder reranker (\texttt{ms-marco-MiniLM-L-6-v2}), and configurable LLM backends for answer generation. Model development is operationalised as comparative system design: baseline vector-only retrieval is contrasted against a hybrid configuration whose retrieval-stage parameters---RRF weights, candidate pool sizing, BM25 quality gate, recency decay, and cross-encoder integration---were tuned through a design-phase grid search and iterative ablation (Section~\ref{sec:sensitivity}). The project therefore addresses machine learning model development through model selection, retrieval-stage tuning, and empirical evaluation under deployment constraints, which is appropriate for a production-oriented RAG system in a cost-constrained academic setting. A conventional supervised training workflow, formal experiment tracker (e.g., MLflow), and model registry were outside the current prototype scope; these are identified as future work in Section~\ref{sec:future}.

\section{Baseline Retrieval Configuration}
\label{sec:baseline}

The baseline configuration performs dense cosine similarity search across the full indexed corpus. Given a query embedding $\mathbf{e}_q$ and document embeddings $\mathbf{e}_i$, cosine similarity is computed as:

\begin{equation}
  \text{Sim}(q, i) = \frac{\mathbf{e}_q \cdot \mathbf{e}_i}{\|\mathbf{e}_q\|\|\mathbf{e}_i\|}
  \label{eq:cosine}
\end{equation}

Top-$N$ documents are selected purely on similarity ranking with a minimum similarity threshold of 0.15 to filter retrieval noise. This low threshold is intentionally permissive in the baseline to avoid artificially handicapping it; borderline documents that would be excluded by a stricter threshold are instead ranked low by subsequent RRF fusion in the hybrid configuration. Under exact search, computational complexity scales as $O(Nd)$, where $N$ is the number of indexed chunks and $d = 384$ (the MiniLM embedding dimension). This configuration serves as the experimental control for retrieval performance comparison.

\section{Implemented Hybrid Retrieval Configuration}
\label{sec:hybrid}

The hybrid configuration introduces five refinement stages applied before generation.

\textbf{Stage 1 --- Metadata Filtering:} Structured SQL constraints reduce the candidate space based on document recency (a 180-day sliding window), source type, and document state. Only documents in the \texttt{INDEXED} state are considered. Users may further restrict retrieval to specific source types via the optional \texttt{sources} parameter on the \texttt{/ask} endpoint (e.g., \texttt{["arxiv",\,"hn"]}); values are validated against a server-side allowlist (\texttt{VALID\_SOURCES}). The 180-day window was chosen to balance corpus recency with sufficient document volume: at the expected ingestion rates across five sources, the window is projected to yield 500--2,000 indexed chunks, providing a meaningful retrieval pool while excluding content more than six months old, which is considered stale for fast-moving AI/ML topics. This pre-filtering stage reduces retrieval noise prior to semantic ranking.

\textit{Robustness to missing and noisy metadata.} The filter degrades gracefully when metadata is incomplete. The recency filter uses an inclusive \texttt{OR} clause (\texttt{published\_at IS NULL OR published\_at >= threshold}), retaining documents with missing dates rather than silently dropping them. Each source adapter normalizes metadata at ingestion (whitespace normalization, URL defaults, timestamp coercion from API-specific fields). In the RRF fusion stage, documents with \texttt{NULL} dates receive age $= 9{,}999$ days, deprioritizing them by recency while preserving their semantic and BM25 retrieval paths.

\textbf{Stage 2 --- Vector Similarity Search:} Cosine similarity (Equation~\ref{eq:cosine}) is applied over the filtered candidate subset using the pgvector \texttt{<=>} operator. To provide sufficient candidates for RRF fusion, $8 \times \text{top\_k}$ results are retrieved at this stage. The wider $8\times$ multiplier (compared to $4\times$ for BM25) gives asymmetric candidate pools: the vector signal retrieves more candidates to ensure adequate overlap with BM25 results during fusion, improving the diversity of the final merged ranking without inflating retrieval latency.

\textbf{Stage 3 --- Parallel BM25 Retrieval:} Concurrently with vector search, a full-corpus BM25 keyword search is executed using the BM25Okapi algorithm (Robertson \& Walker, 1994). Given a query $Q$ containing terms $q_1, q_2, \ldots, q_n$, the BM25 score for document $D$ is:

\begin{equation}
  \text{BM25}(D, Q) = \sum_{i=1}^{n} \text{IDF}(q_i) \cdot \frac{f(q_i, D) \cdot (k_1 + 1)}{f(q_i, D) + k_1 \cdot \bigl(1 - b + b \cdot \frac{|D|}{\text{avgdl}}\bigr)}
  \label{eq:bm25}
\end{equation}

Intuitively, BM25 prioritizes passages that contain query keywords frequently while down-weighting terms that are common across the corpus, providing strong lexical precision for exact or near-exact terminology matches.

\noindent where $f(q_i, D)$ is the term frequency of $q_i$ in $D$, $|D|$ is the document length, $\text{avgdl}$ is the average document length across the corpus, and $k_1 = 1.5$, $b = 0.75$ are the BM25Okapi defaults. Crucially, IDF is computed over the \emph{full} indexed corpus ($N = 6{,}220$ chunks at evaluation time), not over the filtered candidate set---computing IDF over a small candidate pool would produce up to $6\times$ error in term weighting (observed empirically during development). The BM25 index is built once and cached at the module level with a 30-minute TTL, refreshed automatically to capture newly ingested documents.

The same source-type and recency filters applied in Stage~1 are enforced on BM25 candidates in Python, ensuring both retrieval signals operate over an identical filtered corpus.

\textbf{Stage 4 --- Reciprocal Rank Fusion (RRF):} The union of vector and BM25 candidate sets is fused using Reciprocal Rank Fusion (Cormack, Clarke, \& Butt, 2009):

\begin{equation}
  \text{RRF}(d) = \sum_{r \in \mathcal{R}} \frac{1}{K + \text{rank}_r(d)}
  \label{eq:rrf}
\end{equation}

\noindent where $\mathcal{R} = \{\text{vector}, \text{BM25}, \text{recency}\}$ represents three independent rankings, $\text{rank}_r(d)$ is the 1-indexed rank of document $d$ in ranking $r$, and $K = 60$ is the standard RRF constant from the original paper. Recency is incorporated as a third ranking signal by sorting candidates by document age in days (ascending; documents with unknown publication dates receive age $= 9{,}999$ and sort last). Each signal is multiplied by a data-informed weight: $w_{\text{vector}} = 0.50$, $w_{\text{BM25}} = 0.35$, $w_{\text{recency}} = 0.15$, so that the fused score becomes:

\begin{equation}
  \text{RRF}_{\text{weighted}}(d) = \sum_{r \in \mathcal{R}} \frac{w_r}{K + \text{rank}_r(d)}
  \label{eq:weighted_rrf}
\end{equation}

Intuitively, weighted RRF is robust because it fuses multiple ranked lists by position rather than raw score scale, so no single retrieval signal can dominate purely due to score calibration differences.

The asymmetric weights reflect a domain-informed prior: semantic similarity (0.50) dominates because the AI/ML corpus has strong embedding coverage; BM25 (0.35) captures keyword matches missed by semantic embeddings; recency (0.15) is a weak signal to avoid promoting freshly-published but thin content over semantically richer articles. The weights were chosen based on internal ablation experiments conducted during design-phase tuning (Ma et al., 2023; Lin, 2021); detailed results are summarized in Section~\ref{sec:sensitivity}.

\textbf{Stage 5 --- Cross-Encoder Reranking:} After RRF fusion and a post-fusion quality gate (bottom-quartile removal), the top candidates are reranked by a cross-encoder model (\texttt{ms-marco-MiniLM-L-6-v2}; Nogueira \& Cho, 2019). Unlike embedding-based similarity, cross-encoders jointly encode the query and passage, enabling them to capture fine-grained interactions such as negations, exact keyword positions, and partial overlap---distinctions invisible to independent embedding comparison. The cross-encoder is lazy-loaded once and cached for the process lifetime to avoid repeated model loading in serverless environments. This stage adds the final precision layer before URL deduplication and top-$k$ selection.

Figure~\ref{fig:retrieval_comparison} illustrates the three retrieval strategies compared during development.

% ---- TikZ: Three-way retrieval comparison diagram ----
\begin{figure}[H]
\centering
\resizebox{\textwidth}{!}{%
\begin{tikzpicture}[node distance=0.45cm and 1.4cm]

% ---- LEFT: BASELINE ----
\node[pbox=awsgray, minimum width=2.4cm] (bq) {User Query};
\node[pbox=awsblue, below=of bq, minimum width=2.4cm] (bvs) {Vector Similarity\\Search (Full Corpus)};
\node[pbox=awsgreen, below=of bvs, minimum width=2.4cm] (btopk) {Top-$k$ Selection};
\node[pbox=awsorange, below=of btopk, minimum width=2.4cm] (bgen) {LLM\\Generation};

\draw[parrow] (bq) -- (bvs);
\draw[parrow] (bvs) -- (btopk);
\draw[parrow] (btopk) -- (bgen);

\node[layerlabel=awsgray, above=0.3cm of bq, font=\small\bfseries] {\textbf{Baseline (Vector-Only)}};

\begin{scope}[on background layer]
\node[boundbox=awsgray, fit=(bq)(bvs)(btopk)(bgen), inner sep=8pt] {};
\end{scope}

% ---- CENTER: WEIGHTED HYBRID (Design-Phase) ----
\node[pbox=propgray, minimum width=2.4cm, right=4.0cm of bq] (oq) {User Query};
\node[pbox=propgray, below=of oq, minimum width=2.4cm] (omf) {Metadata Filter\\(recency + source)};
\node[pbox=propgray, below=of omf, minimum width=2.4cm] (ovs) {Vector Similarity\\Search (Filtered)};
\node[pbox=propred, below=of ovs, minimum width=2.4cm] (orr) {Weighted\\$\alpha$/$\beta$/$\gamma$ Reranking};
\node[pbox=propgray, below=of orr, minimum width=2.4cm] (otopk) {Top-$k$ Selection};
\node[pbox=propgray, below=of otopk, minimum width=2.4cm] (ogen) {LLM\\Generation};

\draw[parrow, color=propgray!60] (oq) -- (omf);
\draw[parrow, color=propgray!60] (omf) -- (ovs);
\draw[parrow, color=propgray!60] (ovs) -- (orr);
\draw[parrow, color=propgray!60] (orr) -- (otopk);
\draw[parrow, color=propgray!60] (otopk) -- (ogen);

\node[layerlabel=propred, above=0.3cm of oq, font=\small\bfseries] {\textbf{Weighted Hybrid (Design-Phase)}};

\begin{scope}[on background layer]
\node[boundbox=propgray, fit=(oq)(omf)(ovs)(orr)(otopk)(ogen), inner sep=8pt] {};
\end{scope}

% annotation for original hybrid
\draw[-{Stealth}, dashed, color=propred!70, thick]
  ([xshift=0.3cm]orr.east) -- ++(0.4,0)
  node[right, font=\scriptsize, color=propred, text width=1.8cm]
  {p95 = 25\,s\\weights dataset-\\dependent};

% ---- RIGHT: FINAL HYBRID (BM25+RRF) ----
\node[pbox=awsgray, minimum width=2.4cm, right=4.0cm of oq] (hq) {User Query};
\node[pbox=propred, below=of hq, minimum width=2.4cm] (hmf) {Metadata Filter\\(recency + source)};
\node[pbox=awsblue, below=of hmf, minimum width=2.4cm] (hvs) {Vector Similarity\\Search (Filtered)};
\node[pbox=awspurple, below=of hvs, minimum width=2.4cm] (hrr) {BM25 + Weighted\\RRF + Cross-Enc.};
\node[pbox=awsgreen, below=of hrr, minimum width=2.4cm] (htopk) {Top-$k$ Selection};
\node[pbox=awsorange, below=of htopk, minimum width=2.4cm] (hgen) {LLM\\Generation};

\draw[parrow] (hq) -- (hmf);
\draw[parrow] (hmf) -- (hvs);
\draw[parrow] (hvs) -- (hrr);
\draw[parrow] (hrr) -- (htopk);
\draw[parrow] (htopk) -- (hgen);

\node[layerlabel=propblue, above=0.3cm of hq, font=\small\bfseries] {\textbf{Final Hybrid (BM25+RRF)}};

\begin{scope}[on background layer]
\node[boundbox=propblue, fit=(hq)(hmf)(hvs)(hrr)(htopk)(hgen), inner sep=8pt] {};
\end{scope}

% annotation for final hybrid
\draw[-{Stealth}, dashed, color=awspurple!70, thick]
  ([xshift=0.3cm]hrr.east) -- ++(0.4,0)
  node[right, font=\scriptsize, color=awspurple, text width=1.8cm]
  {p95 = 2.44\,s\\weighted RRF\\+ cross-encoder};

\end{tikzpicture}%
}
\caption{Retrieval strategy comparison: baseline (vector-only) vs. hybrid with RRF fusion. The figure illustrates the evolution of retrieval strategies and highlights the latency reduction achieved through RRF-based fusion (p95~2.44\,s) compared to earlier linear weighted reranking approaches (p95~25\,s).}
\label{fig:retrieval_comparison}
\end{figure}

\section{Hallucination Verification Strategy}
\label{sec:prompt}

TechPulse applies a three-layer verification approach to detect and flag potentially hallucinated outputs. The prompt construction and template details are presented in Section~\ref{sec:prompt_ch4}.

\begin{enumerate}[leftmargin=0.5in]
  \item \textbf{Prompt layer:} The system instruction enforces strict evidence grounding and abstention on insufficient context.

  \item \textbf{Evaluation layer:} RAGAS faithfulness scoring decomposes each generated response into atomic claims and verifies whether each claim is supported by at least one retrieved document.

  \item \textbf{Output layer --- Citation traceability:} At runtime, automated citation grounding verification analyzes each response sentence-by-sentence to compute a citation ratio---the fraction of sentences containing valid \texttt{[Source N]} references. Responses falling below the configurable citation grounding threshold are flagged as potentially hallucinated and replaced with an abstention message, with the original response preserved in metadata for review.
\end{enumerate}

\noindent This strategy reduces but does not eliminate hallucination risk, particularly under limited or noisy retrieval contexts where evidence coverage is insufficient.

\section{Optimization Objective}
\label{sec:optimization}

The retrieval task is framed as a constrained optimization problem:

\begin{equation}
  \max \; \text{Precision}(q) \quad \text{subject to} \quad \text{Latency} \leq \tau
  \label{eq:opt}
\end{equation}

where $\tau$ is the p95 latency target of 2 seconds. This formulation makes the precision--latency trade-off explicit and provides a principled basis for tuning retrieval depth and metadata filter selectivity.

\section{Tool Selection and Implementation Stack}
\label{sec:tools}

Table~\ref{tab:tools} summarizes the implemented tools and their justification. The stack has been adapted from the original proposal to operate within AWS Free Tier constraints, as detailed in Chapter~\ref{ch:progress}.

\begin{table}[H]
\caption{Implemented Tools, Services, and Justification.}
\label{tab:tools}
\centering
\begin{tabular}{p{3.2cm}p{3.2cm}p{6.6cm}}
\toprule
\textbf{Component} & \textbf{Tool / Service} & \textbf{Justification} \\
\midrule
Ingestion orchestration & AWS Lambda + EventBridge & Serverless, scheduled/event-driven, no cluster management \\
Message queue & Amazon SQS + DLQ & Decouples ingestion from embedding; failure resilience \\
Raw storage & Amazon S3 (medallion layout) & Cost-effective, durable, prefix-based lineage tracking \\
Embedding model & fastembed (ONNX Runtime; 
\texttt{sentence-transformers/all-MiniLM-L6-v2}) & Low latency, 384-dim, no PyTorch dependency in Lambda container \\
Vector + relational store & RDS PostgreSQL + pgvector & Unified hybrid retrieval; AWS Free Tier compatible \\
Generative model & Groq (Llama 3.1 8B Instant, primary); Bedrock (Amazon Nova Micro, fallback); Ollama / HuggingFace (local fallbacks) & Auto-cascade fallback chain; Groq free-tier API available across all environments \\
Frontend & React (S3 static hosting) & Lightweight, decoupled from backend \\
API routing & Amazon API Gateway (HTTP) & Managed, HTTPS, CORS support \\
CI/CD & GitHub Actions + SAM & Reproducible IaC deployments, automated testing \\
Monitoring & CloudWatch + SNS alarms & Observability, error alerts, health checks \\
Evaluation & RAGAS + statistical testing & Standardized RAG quality metrics + paired comparison \\
\bottomrule
\end{tabular}
\end{table}

\noindent The full LLM backend fallback hierarchy and execution logic are detailed in Section~\ref{sec:llm_fallback} (Chapter~\ref{ch:arch}).

\section{System Context Overview}
\label{sec:system_context}

Figure~\ref{fig:system_context} presents a high-level context view of TechPulse, showing external actors, data flows, and system boundaries before the detailed architecture is described in Chapter~\ref{ch:arch}.

% ---- TikZ: System Context Diagram ----
\begin{figure}[H]
\centering
\begin{tikzpicture}[node distance=1.0cm and 1.6cm, font=\small]

% Central system box
\node[rectangle, draw=awsblue, fill=awslightblue, rounded corners=6pt,
      minimum width=5.5cm, minimum height=3.2cm,
      align=center, font=\small\bfseries, line width=1.2pt]
  (sys) at (0,0) {\textbf{TechPulse}\\
  \scriptsize Hybrid RAG System\\
  \scriptsize (AWS Cloud + Docker Local)};

% External data sources (left) — pushed further out
\node[rectangle, draw=awsgreen!70, fill=awsgreen!15, rounded corners=4pt,
      text width=2.8cm, align=center, minimum height=0.65cm, font=\scriptsize]
  (arxiv) at (-6.2, 1.6) {ArXiv / Hacker News / DEV.to / GitHub / RSS};

% User (right top) — pushed further out
\node[rectangle, draw=awsorange!70, fill=awsorange!15, rounded corners=4pt,
      text width=2.6cm, align=center, minimum height=0.65cm, font=\scriptsize]
  (user) at (6.2, 1.0) {User};

% Evaluation system (right bottom) — pushed further out
\node[rectangle, draw=awspurple!70, fill=awspurple!15, rounded corners=4pt,
      text width=2.6cm, align=center, minimum height=0.65cm, font=\scriptsize]
  (eval) at (6.2, -1.0) {Evaluation (RAGAS + Stats)};

% Cloud / ops (bottom) — pushed further down
\node[rectangle, draw=propgray!70, fill=propgray!15, rounded corners=4pt,
      text width=2.8cm, align=center, minimum height=0.65cm, font=\scriptsize]
  (ops) at (0, -3.2) {Operators (CloudWatch / SNS Alerts)};

% Arrows with white-filled labels for visibility
\draw[-{Stealth}, thick, color=awsgreen!70]
  (arxiv.east) -- node[fill=white, inner sep=2pt, font=\scriptsize\bfseries, text=awsgreen!80] {raw content} (sys.west |- arxiv);

\draw[{Stealth}-{Stealth}, thick, color=awsorange!70]
  (user.west) -- node[fill=white, inner sep=2pt, font=\scriptsize\bfseries, text=awsorange!80, above, pos=0.5] {NL query}
                  node[fill=white, inner sep=2pt, font=\scriptsize\bfseries, text=awsorange!80, below, pos=0.5] {cited response}
  (sys.east |- user);

\draw[-{Stealth}, dashed, thick, color=awspurple!70]
  (sys.east |- eval) -- node[fill=white, inner sep=2pt, font=\scriptsize\bfseries, text=awspurple!80] {eval metrics} (eval.west);

\draw[-{Stealth}, dashed, thick, color=propgray!70]
  (sys.south) -- node[fill=white, inner sep=2pt, font=\scriptsize\bfseries, text=propgray!80, right] {drift alerts / health} (ops.north);

\end{tikzpicture}
\caption{System context diagram: data flows and evaluation signals.}
\label{fig:system_context}
\end{figure}

\medskip
\noindent The following chapters translate these methodological decisions into a concrete system architecture (Chapter~4), evaluation framework and results (Chapter~5), implementation details (Chapter~6), and conclusions (Chapter~7).
